\chapter{Conclusion}
\label{ch:conclusion}

The question given to me at the start of this internship was whether a
fine-tuned language model can predict the next \cid{cell\_id} of a
mobile-network user. Nine weeks and eleven iterations later the answer is yes,
by a modest and measurable margin, and the more interesting part of the answer is
the conditions under which that is true.

\paragraph{What the work settles.} A fine-tuned Qwen2.5-0.5B with LoRA adapters
beats the persistence baseline by \textbf{+4.3 points} of top-1 accuracy,
reproducibly, on two independent held-out test sets and across five training
runs. That number is not large, and it should not be: the competitor is a
one-line rule that is right roughly half the time on mobile users, and for five
of the eleven iterations nothing I trained beat it at all. What made the
difference was neither capacity nor volume but \emph{conditioning on who the
user is}, grouping users by the hours at which they actually change physical
site, and routing each group to its own set of weights. That single change took
the model from $1.8$ points below the rule to $3.7$ points above it, with the
gain distributed in exact inverse proportion to how well repetition already
worked for each profile ($r = -0.99$).

\paragraph{The operating point is small.} Data volume is a real lever over one
narrow stretch, from roughly 400 to 2{,}000 training users, where the edge
grows from $+0.6$ to $+4.2$ points, and worth nothing above it. A further
ten-fold increase, to 21{,}000 users, buys $+0.1\pt$, and a paired bootstrap over
1{,}000 held-out users puts zero inside every confidence interval, including the
one spanning the entire range. One mechanism fits the timing of that transition:
the least populated of the four behavioural profiles crosses roughly a hundred
members exactly between the two tiers where the score moves. Below that, a
profile cannot be estimated and the model is right to ignore it.

\paragraph{Difficulty lives in the user, not in the dataset.} This is the finding
I would keep if I had to keep one. Twenty-two points separate the most sedentary
behavioural profile from the most mobile one, against three tenths of a point
across the whole training-size ladder. Per-user accuracy on this task ranges
from zero to one hundred per cent, and what predicts it is how
repetitive that person's mobility already is. An aggregate accuracy on this
problem is a weighted average over segments of wildly different difficulty, and
quoting it alone hides the only part of the result that is about the model.

\paragraph{And a methodological point that cost more than most of the
modelling.} Three of the largest single improvements of the internship were not
modelling ideas. Reading the best checkpoint instead of the last was worth
$+15.7$ points. Recognising that the context cut is part of the ruler rather than
of the model retired a $+1.1$-point ``result'' that was noise. Choosing a corpus
on which the trivial rule does not already score 69\,\% is what made the question
measurable at all. On a task whose honest yardstick is a one-line heuristic, the
first place to look for points is the measurement.

\paragraph{What it leaves open.} The question given to me was a within-day one,
predict the rest of a person's day from its beginning, and that is the question
this report answers. The obvious continuation is to widen it: predicting day
$N{+}1$ from days $1 \ldots N$, observing a home--activity transition several
times for the same individual, and anything that could be called
personalisation. None of that is a matter of a better model; it needs a corpus
whose identifiers are stable from one day to the next, and it is where I would
take the work next. A seven-day comparison set up to probe part of that question
crashed after five hours without completing a single training, and the report
claims no day-of-week effect as a result.

The honest summary of the nine weeks is therefore narrow and, I think, useful:
on single-day cellular traces, the population prior is cheap, it is measurable,
it is worth about four points over a trivial rule, and it is exhausted at two
thousand users.
